\section{Accuracy, cost and rejection}
\label{sec:accuracy-cost}

\subsection{An analytic model (T6)}
\label{sec:t6}
% Member 6

To see exactly how accuracy, cost and rejection trade off, we use a model in which every
quantity has a closed form. Let
\[
  y \sim \mathrm{Bernoulli}(\pio), \qquad x \mid y=1 \sim \mathcal N(+\mu, 1), \qquad
  x \mid y=0 \sim \mathcal N(-\mu, 1), \qquad \mu > 0 .
\]
Write $\phi$ and $\Phi$ for the standard normal density and CDF, $\piz = 1-\pio$, and
$\lambda = \log(\pio/\piz)$ for the prior log-odds.

\paragraph{Posterior and thresholds on $x$.}
The log-likelihood ratio is
$\log\frac{\phi(x-\mu)}{\phi(x+\mu)} = \tfrac12\big[(x+\mu)^2 - (x-\mu)^2\big] = 2\mu x$, so
\begin{equation}
  \post(x) = \Prob(y=1\mid x) = \sigm(2\mu x + \lambda),
  \qquad
  \frac{\pio\,\phi(x-\mu)}{\piz\,\phi(x+\mu)} = e^{2\mu x+\lambda} = \frac{\post(x)}{1-\post(x)} .
  \label{eq:t6-post}
\end{equation}
Since $\post(x)$ is strictly increasing in $x$, the probability threshold $t$ corresponds to the
threshold
\begin{equation}
  x_t = \frac{\logit t - \lambda}{2\mu}
  \label{eq:t6-xt}
\end{equation}
on $x$ ($x_0=-\infty$, $x_1=+\infty$). Every rule ``deliver if $\post\le\thz$, spam if
$\post\ge\tho$, inspect in between'' is therefore the rule ``deliver if $x\le\tau_0$, spam if
$x\ge\tau_1$, inspect in between'' with $\tau_k = x_{\theta_k}$. A two-action rule is the case
$\tau_0=\tau_1=\tau$.

\paragraph{Closed forms.}
For a rule $(\tau_0,\tau_1)$ the probability of each action given each label is
\begin{center}
\begin{tabular}{lccc}
\toprule
 & deliver & spam folder & inspect \\
\midrule
$y=0$ & $\Phi(\tau_0+\mu)$ & $\Phi(-\tau_1-\mu)$ & $\Phi(\tau_1+\mu)-\Phi(\tau_0+\mu)$ \\
$y=1$ & $\Phi(\tau_0-\mu)$ & $\Phi(\mu-\tau_1)$  & $\Phi(\tau_1-\mu)-\Phi(\tau_0-\mu)$ \\
\bottomrule
\end{tabular}
\end{center}
Weighting the rows by $\piz$ and $\pio$ gives the population confusion matrix, and from it
\begin{align}
  \rej(\tau_0,\tau_1) &= \piz\big[\Phi(\tau_1+\mu)-\Phi(\tau_0+\mu)\big]
                        + \pio\big[\Phi(\tau_1-\mu)-\Phi(\tau_0-\mu)\big], \label{eq:t6-rho}\\
  R(\tau_0,\tau_1)    &= \cfp\,\piz\,\Phi(-\tau_1-\mu) + \cfn\,\pio\,\Phi(\tau_0-\mu) + \ci\,\rej(\tau_0,\tau_1),
                        \label{eq:t6-cost}\\
  \mathrm{Acc}(\tau_0,\tau_1) &= \frac{\piz\,\Phi(\tau_0+\mu) + \pio\,\Phi(\mu-\tau_1)}{1-\rej(\tau_0,\tau_1)} .
                        \label{eq:t6-acc}
\end{align}
Here $\rej$ is the rejection (inspection) rate, $R$ the expected cost per email, and
$\mathrm{Acc}$ the accuracy on the emails that were not inspected. For a two-action rule
$\rej=0$ and $\mathrm{Acc}(\tau) = \piz\Phi(\tau+\mu)+\pio\Phi(\mu-\tau)$.

\emph{Check against Section 2.} Differentiating \eqref{eq:t6-cost} and using \eqref{eq:t6-post},
\begin{align*}
  \frac{\partial R}{\partial \tau_1} &= \piz\phi(\tau_1+\mu)\Big[-(\cfp-\ci) + \ci\, e^{2\mu\tau_1+\lambda}\Big], &
  \frac{\partial R}{\partial \tau_0} &= \piz\phi(\tau_0+\mu)\Big[(\cfn-\ci)\,e^{2\mu\tau_0+\lambda} - \ci\Big].
\end{align*}
Each bracket changes sign once, from negative to positive. So $R$ is minimised where the posterior
odds equal $(\cfp-\ci)/\ci$ and $\ci/(\cfn-\ci)$, that is at $\post=1-\ci/\cfp=\tho$ and
$\post=\ci/\cfn=\thz$. These are the thresholds of Section~2, found here without using the
pointwise argument. (This needs $\thz<\tho$, i.e.\ $\ci<\cstar$; otherwise the optimum is a
two-action rule.)

\subsubsection*{Higher accuracy can mean higher cost}

\begin{proposition}\label{prop:acc-cost}
Let $\tau_{\mathrm{acc}} = x_{1/2} = -\lambda/(2\mu)$ and $\tau_{c} = x_{\tstar}$ with
$\tstar = \cfp/(\cfp+\cfn)$. Among two-action rules, $\tau_{\mathrm{acc}}$ uniquely maximises
accuracy and $\tau_c$ uniquely minimises cost. If $\cfp\neq\cfn$ then $\tau_{\mathrm{acc}}\neq\tau_c$,
and for any two thresholds $\tau<\tau'$ in the closed interval between them, the one nearer
$\tau_{\mathrm{acc}}$ has strictly higher accuracy \emph{and} strictly higher cost. In particular the
accuracy-optimal rule is more accurate but more expensive than the cost-optimal rule.
\end{proposition}

\begin{proof}
Using \eqref{eq:t6-post},
\[
  \mathrm{Acc}'(\tau) = \piz\phi(\tau+\mu)\big[1 - e^{2\mu\tau+\lambda}\big], \qquad
  R'(\tau) = \piz\phi(\tau+\mu)\big[\cfn e^{2\mu\tau+\lambda} - \cfp\big].
\]
The factor $\piz\phi(\tau+\mu)$ is positive and $e^{2\mu\tau+\lambda}$ is strictly increasing. So
$\mathrm{Acc}$ strictly increases up to the root $e^{2\mu\tau+\lambda}=1$ (posterior $\tfrac12$,
$\tau=\tau_{\mathrm{acc}}$) and strictly decreases after it. Likewise $R$ strictly decreases up to
$e^{2\mu\tau+\lambda}=\cfp/\cfn$ (posterior $\tstar$, $\tau=\tau_c$) and strictly increases after it.
If $\cfp\ne\cfn$ then $\tstar\ne\tfrac12$ and, as $x_t$ is strictly increasing, $\tau_c\ne\tau_{\mathrm{acc}}$.
Suppose $\cfp>\cfn$, so $\tau_{\mathrm{acc}}<\tau_c$. On $[\tau_{\mathrm{acc}},\tau_c]$, $\mathrm{Acc}$ is
strictly decreasing and $R$ is strictly decreasing, so moving from $\tau$ towards $\tau_c$ lowers both.
The case $\cfp<\cfn$ is symmetric.
\end{proof}

So between the two optima, accuracy and cost always rank rules in \emph{opposite} order: picking
the more accurate of two rules always picks the more expensive one. With $\mu=1$, $\pio=0.4$,
$\cfp=10$, $\cfn=1$ and $\ci=0.3$, Table~\ref{tab:t6} gives the numbers. The accuracy rule is the
most accurate two-action rule, at $84.6\%$, yet it costs $0.772$ per email. That is more than twice the
cost-optimal rule ($0.311$), and also more than the trivial rule that never sends anything to spam
($60\%$ accuracy, cost $0.400$). Figure~\ref{fig:t6-acc-cost} shows both curves against the
probability threshold.

\begin{table}[ht]
\centering\small
\input{\resultsdir/t6_rule_comparison.tex}
\caption{Analytic model with $\mu=1$, $\pio=0.4$, $\cfp=10$, $\cfn=1$, $\ci=0.3$. Accuracy is on
non-inspected emails; FPR $=\Prob(\text{spam folder}\mid y=0)$, FNR $=\Prob(\text{deliver}\mid y=1)$.
Generated by \texttt{experiments/t6\_analytic.py}.}
\label{tab:t6}
\end{table}

\begin{figure}[ht]
\centering
\includegraphics[width=\linewidth]{t6_accuracy_vs_cost.pdf}
\caption{Accuracy and expected cost of the two-action rule ``spam iff $\post\ge t$'' in the analytic
model. Accuracy peaks at $t=0.5$, while cost is lowest at $\tstar=10/11$.}
\label{fig:t6-acc-cost}
\end{figure}

\subsubsection*{Bayes risk as a function of the inspection cost}

The next result holds for any model with calibrated posteriors, not only the Gaussian one. Let
\[
  m(x) = \min\{\cfn\,\post(x),\ \cfp\,(1-\post(x))\}
\]
be the posterior expected loss of the better automatic action. By Section~2 the Bayes action
inspects exactly when $\ci < m(x)$, so the Bayes risk and the rejection rate are
\begin{equation}
  \Rstar(\ci) = \E\big[\min\{m(x),\,\ci\}\big], \qquad \rej(\ci) = \Prob\big(m(x) > \ci\big).
  \label{eq:t6-bayes-risk}
\end{equation}
Write $\Rtwo = \E[m(x)]$ for the Bayes risk without inspection and
$\cstar = \cfp\cfn/(\cfp+\cfn)$.

\begin{lemma}\label{lem:mmax}
$m(x)\le\cstar$ for every $x$, with equality iff $\post(x)=\tstar$.
\end{lemma}
\begin{proof}
$\cfn p$ increases and $\cfp(1-p)$ decreases in $p$. Their minimum is therefore largest where
they are equal, at $p=\cfp/(\cfp+\cfn)=\tstar$, where both equal $\cstar$.
\end{proof}

\begin{proposition}\label{prop:risk-ci}
\begin{enumerate}[label=(\roman*),itemsep=1pt]
  \item $\Rstar$ is non-decreasing and concave on $[0,\infty)$, with $\Rstar(0)=0$.
  \item Its right derivative is $\rej(\ci)$. Wherever $\Prob(m(x)=\ci)=0$ (for every $\ci$ in the Gaussian
        model), $\Rstar$ is differentiable with $\frac{d\Rstar}{d\ci} = \rej(\ci)$, and $\rej$ is
        non-increasing.
  \item Rejection stops at $\cstar$: $\rej(\ci)=0$ and $\Rstar(\ci)=\Rtwo$ for all $\ci\ge\cstar$.
  \item $\Rtwo-\Rstar(\ci)=\int_{\ci}^{\cstar}\rej(s)\,ds$. In the Gaussian model this is strictly
        positive iff $\ci<\cstar$, so allowing inspection improves the system iff $\ci<\cstar$.
\end{enumerate}
\end{proposition}

\begin{proof}
(i) For each fixed $x$, $c\mapsto\min\{m(x),c\}$ is the minimum of two non-decreasing affine
functions of $c$, so it is non-decreasing and concave. Both properties survive taking the expectation.
At $c=0$ the minimum is $0$ because $m\ge0$.

(ii) For $h>0$,
$\frac1h\big[\min\{m,c+h\}-\min\{m,c\}\big]$ lies in $[0,1]$ and tends to $\mathbf 1\{m>c\}$ as
$h\downarrow0$. By dominated convergence the right derivative of $\Rstar$ is $\Prob(m>c)=\rej(c)$.
The left derivative is, in the same way, $\Prob(m\ge c)$, and the two agree when $\Prob(m=c)=0$.
In the Gaussian model $\post(x)$ is strictly monotone in $x$, so $\{x: m(x)=c\}$ has at most
two points and probability zero. $\rej$ is non-increasing because $\{m>c'\}\subseteq\{m>c\}$ for
$c'>c$. This is a second proof of concavity.

(iii) By Lemma~\ref{lem:mmax}, $m(x)\le\cstar\le\ci$, so $\rej(\ci)=\Prob(m>\ci)=0$ and
$\min\{m,\ci\}=m$, giving $\Rstar(\ci)=\E[m]=\Rtwo$.

(iv) By (ii) and (iii), $\Rtwo-\Rstar(\ci)=\Rstar(\cstar)-\Rstar(\ci)=\int_{\ci}^{\cstar}\rej(s)\,ds$.
In the Gaussian model $\post(x)$ takes every value in $(0,1)$. For $s<\cstar$ the set
$\{x: m(x)>s\}=\{x:\thz(s)<\post(x)<\tho(s)\}$ is a non-empty open interval of $x$, so
$\rej(s)>0$. The integral is therefore positive iff $\ci<\cstar$.
\end{proof}

Proposition~\ref{prop:risk-ci} says three things. First, the rejection rate is the marginal cost
of inspection: raising $\ci$ by a small $\delta$ raises the Bayes risk by $\rej\,\delta$, because only the
inspected emails pay the extra cost. Second, the saving from allowing inspection is the area under
the rejection curve to the right of $\ci$. Third, inspection is worth paying for exactly when it is
cheaper than the worst-case automatic loss $\cstar$. With $\cfp=10$ and $\cfn=1$ that cut-off
is $\cstar=10/11\approx0.909$, and it does not depend on $\mu$ or $\pio$.
Because these results hold for every $\ci$, the price $\ci=0.3$ of Section~\ref{sec:loss-matrix} serves
only as a representative value for tables: it lies well below $\cstar$, so inspection is used, and any
other price below $\cstar$ leads to the same conclusions.
In the Gaussian model the curves follow in closed form: put $\tau_k=x_{\theta_k(\ci)}$ from
\eqref{eq:t6-xt} into \eqref{eq:t6-rho} and \eqref{eq:t6-cost}. Figure~\ref{fig:t6-risk} plots them for
three class separations. Better-separated classes have lower risk at every $\ci$, and all three
curves stop rejecting at the same $\cstar$. The rejection rate, however, is not ordered by $\mu$:
weakly separated classes ($\mu=0.5$) inspect more emails than $\mu=1$ when inspection is cheap, but
fewer once $\ci\gtrsim0.57$. With weak separation most posteriors stay close to the prior $0.4$.
Once the lower threshold $\thz=\ci/\cfn$ climbs past them, those emails are delivered rather than
inspected. (For $\mu=0.5$ the best two-action rule is almost ``always deliver'': $\Rtwo=0.399$
against $\pio\cfn=0.4$.) At the default $\ci=0.3$
with $\mu=1$, inspecting $41.6\%$ of emails cuts the cost from $\Rtwo=0.311$ to $0.179$.

\begin{figure}[ht]
\centering
\includegraphics[width=\linewidth]{t6_risk_reject_vs_cinsp.pdf}
\caption{Analytic Bayes risk $\Rstar(\ci)$ and rejection rate $\rej(\ci)$ for $\pio=0.4$,
$\cfp=10$, $\cfn=1$ and three separations $\mu$. Dotted lines show $\Rtwo$, the best risk without
inspection. Rejection falls to zero at $\cstar=10/11$ (dashed), after which $\Rstar=\Rtwo$.}
\label{fig:t6-risk}
\end{figure}

\begin{assumptions}
One-dimensional equal-variance Gaussian classes with known $\mu$ and $\pio$, so the posterior is
exact (calibrated); zero loss for correct decisions; constant inspection cost $\ci$ and a perfect
inspector. Lemma~\ref{lem:mmax} and Proposition~\ref{prop:risk-ci}(i)--(iii) need only calibrated
posteriors and this loss matrix. Part (iv) and Proposition~\ref{prop:acc-cost} use the Gaussian model.
\end{assumptions}

\emph{Numerical check.} \texttt{tests/test\_analytic.py} draws $400{,}000$ samples from the model,
applies the Bayes rule, and confirms that the realized loss, rejection rate and accuracy match
\eqref{eq:t6-rho}--\eqref{eq:t6-acc}. It also checks that a grid search over all rules
$(\tau_0,\tau_1)$ never beats the Bayes thresholds, and that a finite-difference slope of
$\Rstar$ equals $\rej$.

\subsection{Experiments (C6)}
\label{sec:c6}
For split $s=0,\ldots,19$ we drew $10{,}000$ Gaussian messages ($\mu=1$, $\pio=0.4$,
seed $10{,}000+s$) and split the Spambase corpus using the same stratified $60/20/20$ protocol.
We standardized features and fitted models on training data. Validation log-loss selected
the shared-covariance Gaussian shrinkage from the team's grid and the logistic penalty from
$0.01,0.1,1,10,100$. Each fitted model used identical test predictions across rules and prices.
Tables and bands give means and one standard error over splits. Overlapping Spambase test sets
make these bands descriptive of split variability.

\begin{table}[htbp]
\centering\small
\input{\resultsdir/c6_rules.tex}
\caption{Two automatic rules on held-out messages. Accuracy is the fraction of all test
messages classified correctly. Cost is realized loss per message at $\cfp=10,\cfn=1$.}
\label{tab:c6-rules}
\end{table}

On synthetic data, changing the Gaussian rule from $0.5$ to $10/11$ lowered accuracy
from $84.6\%$ to $73.9\%$ and more than halved its cost (Table~\ref{tab:c6-rules}).
At ten times the price of a missed spam email, avoiding one false positive pays for up to
ten extra false negatives. Both fitted models show the same accuracy--cost ordering
on Spambase.

\begin{table}[htbp]
\centering\small
\input{\resultsdir/c6_savings.tex}
\caption{Inspection at $\ci=0.3$. The two-action baseline uses each fitted posterior at
$\tstar=10/11$. Costs and paired savings are per email. The last column is the fraction inspected.}
\label{tab:c6-savings}
\end{table}

\begin{figure}[htbp]
\centering
\includegraphics[width=0.88\linewidth]{c6_inspection_synthetic.pdf}
\caption{Synthetic held-out results against the exact risk and rejection rate for the
true Gaussian posterior. Narrow bands and overlapping lines reflect close agreement.
The vertical line marks $\cstar=10/11$.}
\label{fig:c6-synthetic}
\end{figure}
\begin{figure}[htbp]
\centering
\includegraphics[width=0.88\linewidth]{c6_inspection_spambase.pdf}
\caption{Spambase held-out cost and inspection rate. At and above the marked cutoff both
models take only automatic actions, so each cost curve levels off at its own baseline.}
\label{fig:c6-spambase}
\end{figure}

At $\ci=0.3$ theory predicts cost $0.179$ and rejection $0.416$. The fitted Gaussian
gives $0.181$ and $0.417$ (Figure~\ref{fig:c6-synthetic}). On Spambase the logistic fit
saves $0.100$ per email while inspecting $24.1\%$ (Figure~\ref{fig:c6-spambase} and Table~\ref{tab:c6-savings}).
At zero price, both models inspected every test email on both datasets, with zero realized
cost in every split. A finite fitted log-odds score implies a probability strictly between
zero and one. Comparing scores directly with the log-odds thresholds preserves this fact
even for forecasts that would round to an endpoint in probability space.

For a true posterior the gain over the best two-action rule is
$\Rtwo-\Rstar(\ci)=\E[(m(X)-\ci)_+]$, where $u_+=\max(u,0)$.
It is strictly positive exactly when $\Prob(m(X)>\ci)>0$. Every nonempty inspection
band has positive probability in the Gaussian model. On Spambase the true posterior is unknown,
so the table compares the fitted policies on observed labels.

\paragraph{Sensitivity to repeated feature vectors.}
About $9.8\%$ of Spambase test rows had an exact 57-feature match in their training split.
Excluding those rows preserved the conclusion: inspection reduced logistic cost from
$0.264$ to $0.160$ and Gaussian cost from $0.291$ to $0.173$ per email.
